%! AP Statistics — Unit 3, Lesson 3.3 — Warm-Up
\documentclass[10pt]{article}
\usepackage{apstats-article}
\usepackage{apstats-boxes}

\begin{document}

\pageheader{Unit 3, Lesson 3.3}{Warm-Up}
\namedateperiod

\vspace{0.3in}

% ── SECTION 1: Spiral Review ──────────────────────────────────────────────────
{\large\bfseries\color{navy} Part 1 \quad Spiral Review}

\vspace{0.12in}

\begin{enumerate}

    \item From Lesson 3.2 --- for each study below, identify the type (observational
    or experiment), and state whether a causal conclusion is justified.

    \vspace{8pt}
    \renewcommand{\arraystretch}{1.8}
    \begin{tabularx}{\linewidth}{@{} X p{2.8cm} p{2.4cm} @{}}
    \textbf{Study} & \textbf{Type} & \textbf{Causal?} \\
    \hline
    Researchers randomly assign 80 students to two test-prep programs and compare scores. &
    \blank{2.5cm} & Yes \quad No \\
    A researcher surveys 500 adults about their diet and health outcomes. &
    \blank{2.5cm} & Yes \quad No \\
    \end{tabularx}

    \vspace{0.18in}

    \item A researcher wants to survey 100 students from a school of 800. She assigns
    each student a number from 1 to 800, then uses a calculator to randomly select
    100 numbers.
    \begin{enumerate}[label=(\alph*), itemsep=4pt]
        \item What sampling method is this? \blank{5cm}
        \item What is one advantage of this method?\\
        \writeline
    \end{enumerate}

\end{enumerate}

\vspace{0.35in}

% ── SECTION 2: Looking Ahead ──────────────────────────────────────────────────
{\large\bfseries\color{navy} Part 2 \quad Looking Ahead}

\vspace{0.12in}

\noindent Today we develop four random sampling methods.

\vspace{0.14in}

\begin{tcolorbox}[colback=greenbg, colframe=greenacc, boxrule=0.8pt, arc=2mm,
    left=4mm, right=4mm, top=3mm, bottom=3mm]
    \small
    \begin{itemize}[leftmargin=1.3em, itemsep=8pt]
  \item \textbf{SRS}: every group of size $n$ has an equal chance of being chosen.
  \item \textbf{Stratified}: divide into homogeneous strata; SRS within each.
  \item \textbf{Cluster}: divide into heterogeneous clusters; randomly select clusters.
  \item \textbf{Systematic}: random start, then every $k$th individual.
    \end{itemize}
\end{tcolorbox}

\vspace{0.2in}

\noindent\textcolor{slate}{\small What is one question you have about choosing between
these methods?}\\[8pt]
\writeline\\[2pt]\writeline

\end{document}
